\chapter*{Preface}
\addcontentsline{toc}{chapter}{Preface}

This book is for ecologists, taxonomists, biodiversity informaticians and the
people who keep the databases running. It assumes you can write a script and
that someone has recently asked you what you are doing about AI. It does not
assume you want to become a machine learning researcher.

This is a pragmatic primer on how the work of ecologists changes as language
models and coding agents enter it. It asks how the tools, the landscape and
the challenges are changing, and what remains at the core of the work. In its
current form it is mostly about ecological informatics, with the literature as
the main source of data. Image, sound and species distribution models appear
only briefly (\Cref{ch:applications}), and molecular and microbial biodiversity
are outside its scope beyond that chapter. For image and sound,
\textcite{tuia2022perspectives} describe machine learning for camera traps,
acoustic recorders and other sensors in animal ecology.
\textcite{pichler2023review} review machine and deep learning for ecologists
more broadly, species distribution models included.

\section*{Three arguments}

\textbf{The real change is not the chatbot.} It is the loop that lets a model
read files, run commands, see the result and try again. Your job moves from
writing code to specifying it, limiting what it can touch, and checking what
comes back --- closer to running a system than to programming.

\textbf{Data is the bottleneck, and ecology has already done much of the work.}
``AI-ready data'' arrived late to a community that has spent thirty years
building Darwin Core, GBIF, ontologies and curated databases by hand. It is worth
asking who built them, and on what terms they are now used. The formats and
vocabularies exist; whether a dataset records what a model needs, such as
sampling effort and non-detections, is a separate question
(\Cref{sec:fitformodel}).

\textbf{Running models on your own machine is a reasonable choice.} It costs
more up front and less per query, keeps unpublished data at home, survives a
service being withdrawn, and works on a research vessel with no connection.

The book does not argue that ecology should adopt AI, nor that it should
refuse. One of the papers behind it argues that ecology is \emph{not} ready
\parencite{mishra2026notready}, and that argument is taken seriously throughout.

\section*{How to read it}

Part~I covers what these systems are, why they arrived now, what they do well
and where they fail. Part~II is data. Part~III is working with agents. Part~IV
is running models locally. Part~V covers security, integrity and keeping up,
and an epilogue sets the book in a wider frame (\Cref{ch:epilogue}).
Chapters are short and mostly independent. \Cref{ch:firstsession,ch:testing}
are exercises with invented data, and \Cref{ch:weeklyliterature} builds a
weekly literature review from real records; the R scripts for all three are in
the book's \code{workshop} folder. If you read only one chapter, read
\Cref{ch:powers}. Model names and prices were correct when written and will
date; the reasoning should not.

\section*{How this book was written}

I used AI agents to draft text, read the literature and assemble the
bibliography, following the practices described in this book. Every reference
corresponds to a document in the project's literature folder; nothing was cited
from a model's memory. A script resolves every DOI through Crossref, or
DataCite where Crossref has no record, and every arXiv identifier through the
arXiv API. It checks that the title, the first author and the year returned
match the entry. References with neither identifier were checked by hand, and
I checked each claim against the source cited for it. The errors that remain
are mine.

\section*{Acknowledgements}

The reading behind this was gathered over two years of the Kybele project and
ELIXIR-GR work at HCMR.
